\chapter{Introduction}
\section{Problem Statement}

Organizations now handle a large number of documents that contain personal data: identity cards, application forms, bank records, medical paperwork, scanned PDFs, and photographs of documents captured on phones. These files often include Aadhaar numbers, PAN numbers, payment card numbers, email addresses, phone numbers, names, and addresses. When such documents are shared without proper redaction, the exposed data can be used for identity theft, fraud, or unauthorized profiling.

Manual redaction is still common, but it does not scale well. A reviewer has to inspect every page, identify sensitive fields, and apply masking without missing a digit or masking the wrong information. The work is repetitive, and mistakes are hard to notice after the document has already been shared. The risk is higher for scanned documents because the sensitive value may appear at an unusual angle, inside a noisy image, or in a layout that differs from the expected template.

The regulatory pressure around privacy has also increased. The GDPR has led to large penalties in Europe, including the EUR 1.2 billion penalty issued against Meta Ireland in May 2023 for unlawful transfers of EU/EEA user data. In India, the Digital Personal Data Protection Act, 2023 sets duties for data fiduciaries and provides monetary penalties that may extend up to Rs. 250 crore for certain breaches. These rules make reliable document handling a practical compliance need, not just a technical preference.

This project focuses on automated redaction for Indian identity-document workflows. Aadhaar is a 12-digit identifier issued to Indian residents. PAN is a ten-character alphanumeric identifier issued by the Income Tax Department, with a fixed structure and a fourth character that encodes holder status. Payment card numbers follow internationally used numbering structures and can be checked with the Luhn algorithm. These identifiers appear frequently in forms and scanned records, so they are good targets for a validation-based redaction system.

Many existing tools support general redaction, but they usually need manual selection, cloud processing, or custom recognizers for Indian identifiers. Generic pattern matching can also over-redact ordinary numbers because it cannot tell whether a 12-digit sequence is a valid Aadhaar number. RE-DACT addresses this gap by combining OCR with validation checks, orientation handling, configurable masking, and an audit log.

\section{Motivation}

The motivation for RE-DACT comes from a simple problem: sensitive Indian identity data is easy to collect and hard to handle safely once it appears inside documents. Government offices, educational institutions, banks, clinics, and private businesses may all receive identity proofs during onboarding, verification, claims, admissions, or compliance checks. A small error during sharing can expose data that should have been hidden.

The project also shows how established tools can be combined into a practical system. Tesseract provides OCR and character-level bounding boxes. OpenCV supports image rotation and masking. Verhoeff and Luhn validation add deterministic checks for structured numbers. FastAPI makes the processing pipeline available through a browser interface and REST endpoints. None of these pieces is new by itself; the value is in connecting them in a way that works for the target document types.

Document orientation is another reason for building the system. Real inputs are rarely perfect. A user may upload a phone photograph, a sideways scan, or a PDF made from images with inconsistent page rotation. A redaction tool that assumes upright text can miss sensitive values. RE-DACT therefore tries multiple orientations and uses an early-exit strategy so normal inputs do not pay the cost of all OCR attempts.

The project is also useful as a learning exercise in applied AI and privacy engineering. It brings together OCR, deterministic validation, web APIs, frontend design, temporary file handling, and hash-based audit logging. The result is small enough to understand, but complete enough to show the trade-offs involved in real document processing.

\section{Project Objectives}

The main objective of this project is to design and implement a web-based platform that detects and masks sensitive information in PDF and image documents, with a focus on Indian identity data.

The specific objectives are:

\begin{enumerate}
    \item Build a PII detection module for Aadhaar, PAN, payment card numbers, email addresses, phone numbers, and optional NER-based names and locations.
    \item Validate Aadhaar candidates with the Verhoeff checksum and payment card candidates with the Luhn checksum to reduce false positives.
    \item Detect PAN numbers through structural pattern matching and fourth-character holder-type validation.
    \item Process PDF, JPG, JPEG, and PNG files through a pipeline that preserves page structure where possible.
    \item Handle rotated inputs by trying four orientations and two preprocessing states.
    \item Provide three masking levels: standard, partial, and full.
    \item Expose the system through REST API endpoints and a browser interface.
    \item Record processing events in a tamper-evident audit log using SHA-256 hash chaining.
    \item Measure detection behavior and processing time on representative sample documents.
\end{enumerate}

\section{Scope of Project}

The project includes OCR-based redaction for structured Indian identifiers and a limited set of general PII types. The implemented system processes PDF, JPG, JPEG, and PNG files. PDF pages are converted to images at 300 DPI, processed page by page, and assembled back into a PDF when redaction is applied. Image inputs are processed directly.

The detection scope includes Aadhaar numbers, payment card numbers, PAN numbers, email addresses, and Indian phone-number-like patterns. When the optional transformer dependencies are installed, the system can also attempt NER-based detection of names and locations through the \texttt{dslim/bert-base-NER} model. This NER support is useful but should be treated as probabilistic and dependent on OCR quality.

The masking scope includes three levels. Standard masking hides all but the last four characters. Partial masking hides the first four characters only, which is the behavior implemented in the current code and user interface. Full masking hides the entire detected value. For names, locations, email addresses, and phone numbers, the implementation masks the full value regardless of the selected level.

The system stores uploaded and processed files only in temporary working directories. Redacted results remain available for download for ten minutes and are then cleaned up. This is a short-retention design rather than permanent storage. The audit log stores metadata such as filename, file type, detection counts, processing time, and hash-chain fields; it does not store the document contents.

The project does not currently support video redaction, handwritten text recognition, Word documents, multi-page TIFF files, scheduled batch jobs, or distributed processing. It also does not verify whether a PAN or Aadhaar number is actually issued to a person. The checks confirm format and checksum validity where applicable, not identity authenticity.

The implementation assumes that the input contains reasonably clear printed or typed text. Low-resolution images, glare, stylized fonts, strong background graphics, and poor contrast can reduce OCR quality and therefore reduce detection accuracy. The system is suitable for an academic prototype and moderate local use; high-volume production deployment would need explicit file-size limits, rate limiting, persistent job storage, stronger audit-chain recovery across restarts, and more robust operational monitoring.

\section{Organization of the Report}

This report is organized into five chapters.

Chapter 1 introduces the problem, motivation, objectives, and scope of RE-DACT.

Chapter 2 reviews related work in document redaction, OCR, PII detection, validation algorithms, and privacy requirements.

Chapter 3 explains the system design and methodology, including the OCR pipeline, validation logic, masking behavior, overlap handling, audit logging, and API design.

Chapter 4 describes the implementation and presents test results, processing-time observations, interface behavior, and comparison with related tools.

Chapter 5 concludes the report and outlines future improvements such as stronger NER support, multilingual OCR, video redaction, batch processing, and containerized deployment.

The References section lists the sources used for the technical and regulatory background. The appendices contain institutional material such as paper and plagiarism-report sections where required.
